\chapter{Lexical Structure}
\label{ch:lexical}

This chapter defines the lexical grammar of \haira{}: the rules for converting source text into tokens.

\section{Source Encoding}

\haira{} source files are encoded in UTF-8. The file extension is \code{.haira}.

\begin{lstlisting}[language=bash,numbers=none]
hello.haira
utils/string_helpers.haira
\end{lstlisting}

\section{Line Structure}

\haira{} is a line-oriented language. Statements are terminated by newlines, not semicolons.

\begin{lstlisting}
x = 1
y = 2
z = x + y
\end{lstlisting}

\subsection{Line Continuation}

Lines can be continued with a trailing backslash or when the line ends with an operator or open bracket:

\begin{lstlisting}
// Explicit continuation
long_result = very_long_function_name(arg1, arg2) \
    + another_function(arg3)

// Implicit continuation (ends with operator)
total = first_value +
    second_value +
    third_value

// Implicit continuation (open bracket)
users = [
    "alice",
    "bob",
    "charlie"
]
\end{lstlisting}

\section{Comments}

\subsection{Single-Line Comments}

\begin{lstlisting}
// This is a single-line comment
x = 42  // Inline comment
\end{lstlisting}

\subsection{Multi-Line Comments}

\begin{lstlisting}
/*
    This is a multi-line comment.
    It can span multiple lines.
*/
\end{lstlisting}

Multi-line comments can be nested:

\begin{lstlisting}
/*
    Outer comment
    /* Nested comment */
    Still in outer comment
*/
\end{lstlisting}

\subsection{Documentation Comments}

\begin{lstlisting}
/// This is a documentation comment for the following item.
/// It can span multiple lines.
fn calculate_area(width: float, height: float) -> float {
    return width * height
}
\end{lstlisting}

\section{Tokens}

The lexer produces the following token types:

\begin{enumerate}
    \item Keywords
    \item Identifiers
    \item Literals (numbers, strings, booleans)
    \item Operators
    \item Delimiters
\end{enumerate}

\section{Keywords}

The following identifiers are reserved as keywords:

\begin{table}[h]
\centering
\begin{tabular}{lllll}
\toprule
\multicolumn{5}{l}{\textit{General-purpose keywords}} \\
\midrule
\code{and} & \code{as} & \code{break} & \code{catch} & \code{chan} \\
\code{continue} & \code{defer} & \code{else} & \code{enum} & \code{false} \\
\code{fn} & \code{for} & \code{from} & \code{if} & \code{import} \\
\code{in} & \code{match} & \code{nil} & \code{not} & \code{or} \\
\code{return} & \code{select} & \code{spawn} & \code{struct} & \code{true} \\
\code{try} & \code{type} & \code{while} & & \\
\midrule
\multicolumn{5}{l}{\textit{Agentic keywords}} \\
\midrule
\code{agent} & \code{provider} & \code{tool} & \code{workflow} & \\
\bottomrule
\end{tabular}
\caption{Reserved keywords}
\end{table}

\begin{note}
The four agentic keywords (\keyword{agent}, \keyword{provider}, \keyword{tool}, \keyword{workflow}) are first-class language constructs for building AI-powered systems. See Chapters~\ref{ch:providers}--\ref{ch:workflows}.
\end{note}

\section{Identifiers}

Identifiers name variables, functions, types, and modules.

\begin{lstlisting}[language=EBNF]
identifier     = letter { letter | digit | "_" } ;
letter         = "a".."z" | "A".."Z" | "_" ;
digit          = "0".."9" ;
\end{lstlisting}

\textbf{Valid identifiers:}
\begin{lstlisting}[numbers=none]
x
name
userName
user_name
_private
Point3D
MAX_SIZE
\end{lstlisting}

\textbf{Invalid identifiers:}
\begin{lstlisting}[numbers=none]
3d_point    // Cannot start with digit
my-name     // Hyphens not allowed
class       // Reserved keyword (if added)
\end{lstlisting}

\subsection{Naming Conventions}

\begin{table}[h]
\centering
\begin{tabular}{ll}
\toprule
\textbf{Entity} & \textbf{Convention} \\
\midrule
Variables & \code{snake\_case} \\
Functions & \code{snake\_case} \\
Types/Structs & \code{PascalCase} \\
Constants & \code{SCREAMING\_SNAKE\_CASE} \\
Modules & \code{snake\_case} \\
\bottomrule
\end{tabular}
\caption{Naming conventions}
\end{table}

\section{Literals}

\subsection{Integer Literals}

\begin{lstlisting}[language=EBNF]
integer_literal = decimal_literal
                | hex_literal
                | octal_literal
                | binary_literal ;

decimal_literal = digit { digit | "_" } ;
hex_literal     = "0" ("x"|"X") hex_digit { hex_digit | "_" } ;
octal_literal   = "0" ("o"|"O") octal_digit { octal_digit | "_" } ;
binary_literal  = "0" ("b"|"B") binary_digit { binary_digit | "_" } ;

hex_digit       = digit | "a".."f" | "A".."F" ;
octal_digit     = "0".."7" ;
binary_digit    = "0" | "1" ;
\end{lstlisting}

\textbf{Examples:}
\begin{lstlisting}
decimal = 42
with_underscores = 1_000_000
hex = 0xFF
octal = 0o755
binary = 0b1010_1100
\end{lstlisting}

\begin{note}
All integer literals are of type \type{int} (64-bit signed).
\end{note}

\subsection{Floating-Point Literals}

\begin{lstlisting}[language=EBNF]
float_literal = decimal_literal "." decimal_literal [ exponent ]
              | decimal_literal exponent ;

exponent = ("e"|"E") ["+"|"-"] decimal_literal ;
\end{lstlisting}

\textbf{Examples:}
\begin{lstlisting}
pi = 3.14159
scientific = 6.022e23
negative_exp = 1.6e-19
\end{lstlisting}

\subsection{Boolean Literals}

\begin{lstlisting}
is_valid = true
is_empty = false
\end{lstlisting}

\subsection{String Literals}

\begin{lstlisting}[language=EBNF]
string_literal = '"' { string_char } '"' ;
string_char    = any_char - ('"' | '\' | newline)
               | escape_sequence ;
\end{lstlisting}

\textbf{Escape sequences:}

\begin{table}[h]
\centering
\begin{tabular}{ll}
\toprule
\textbf{Sequence} & \textbf{Meaning} \\
\midrule
\code{\textbackslash n} & Newline (LF) \\
\code{\textbackslash r} & Carriage return (CR) \\
\code{\textbackslash t} & Horizontal tab \\
\code{\textbackslash\textbackslash} & Backslash \\
\code{\textbackslash"} & Double quote \\
\code{\textbackslash 0} & Null character \\
\code{\textbackslash xNN} & Hex byte (00-FF) \\
\code{\textbackslash u\{NNNN\}} & Unicode code point \\
\bottomrule
\end{tabular}
\caption{String escape sequences}
\end{table}

\textbf{Examples:}
\begin{lstlisting}
simple = "Hello, World!"
with_escapes = "Line 1\nLine 2"
unicode = "Hello \u{1F44B}"  // Wave emoji
\end{lstlisting}

\subsection{String Interpolation}

Strings can contain interpolated expressions using \code{\$\{...\}}:

\begin{lstlisting}
name = "Alice"
age = 30
greeting = "Hello, ${name}! You are ${age} years old."
\end{lstlisting}

\begin{implementation}
String interpolation is desugared to concatenation during parsing:
\begin{lstlisting}
greeting = "Hello, " + name + "! You are " + to_string(age) + " years old."
\end{lstlisting}
\end{implementation}

\subsection{Raw Strings}

Raw strings preserve backslashes literally:

\begin{lstlisting}
regex_pattern = r"^\d{3}-\d{4}$"
windows_path = r"C:\Users\Alice\Documents"
\end{lstlisting}

\subsection{Triple-Quoted Strings}

Triple-quoted strings use \code{"""..."""} for multi-line content:

\begin{lstlisting}
query = """
    SELECT *
    FROM users
    WHERE active = true
"""
\end{lstlisting}

The leading whitespace common to all lines is stripped.

Triple-quoted strings are also used for:
\begin{itemize}
    \item Tool descriptions (mandatory, see Chapter~\ref{ch:tools})
    \item Agent system prompts (optional, for long prompts, see Chapter~\ref{ch:agents})
\end{itemize}

\begin{lstlisting}
tool search(query: string) -> [Result] {
    """Search the web for information"""
    // ...
}

agent Assistant {
    model: anthropic
    system: """
        You are a helpful assistant.
        Always be polite and concise.
    """
}
\end{lstlisting}

\subsection{Nil Literal}

\begin{lstlisting}
empty: User? = nil
\end{lstlisting}

\section{Operators}

\subsection{Arithmetic Operators}

\begin{table}[h]
\centering
\begin{tabular}{ll}
\toprule
\textbf{Operator} & \textbf{Description} \\
\midrule
\code{+} & Addition \\
\code{-} & Subtraction (or unary negation) \\
\code{*} & Multiplication \\
\code{/} & Division \\
\code{\%} & Modulo (remainder) \\
\bottomrule
\end{tabular}
\end{table}

\subsection{Comparison Operators}

\begin{table}[h]
\centering
\begin{tabular}{ll}
\toprule
\textbf{Operator} & \textbf{Description} \\
\midrule
\code{==} & Equal \\
\code{!=} & Not equal \\
\code{<} & Less than \\
\code{>} & Greater than \\
\code{<=} & Less than or equal \\
\code{>=} & Greater than or equal \\
\bottomrule
\end{tabular}
\end{table}

\subsection{Logical Operators}

\begin{table}[h]
\centering
\begin{tabular}{ll}
\toprule
\textbf{Operator} & \textbf{Description} \\
\midrule
\code{and}, \code{\&\&} & Logical AND \\
\code{or}, \code{||} & Logical OR \\
\code{not}, \code{!} & Logical NOT \\
\bottomrule
\end{tabular}
\end{table}

\subsection{Assignment Operators}

\begin{table}[h]
\centering
\begin{tabular}{ll}
\toprule
\textbf{Operator} & \textbf{Description} \\
\midrule
\code{=} & Assignment \\
\code{+=} & Add and assign \\
\code{-=} & Subtract and assign \\
\code{*=} & Multiply and assign \\
\code{/=} & Divide and assign \\
\bottomrule
\end{tabular}
\end{table}

\subsection{Bitwise Operators}

\begin{table}[h]
\centering
\begin{tabular}{ll}
\toprule
\textbf{Operator} & \textbf{Description} \\
\midrule
\code{\&} & Bitwise AND \\
\code{|} & Bitwise OR \\
\code{\^{}} & Bitwise XOR \\
\code{\~{}} & Bitwise NOT (complement) \\
\code{<<} & Left shift \\
\code{>>} & Right shift (arithmetic for signed, logical for unsigned) \\
\code{>>>} & Logical right shift (always fills with zeros) \\
\bottomrule
\end{tabular}
\end{table}



\subsection{Bitwise Assignment Operators}

\begin{table}[h]
\centering
\begin{tabular}{ll}
\toprule
\textbf{Operator} & \textbf{Description} \\
\midrule
\code{\&=} & Bitwise AND and assign \\
\code{|=} & Bitwise OR and assign \\
\code{\^{}=} & Bitwise XOR and assign \\
\code{<<=} & Left shift and assign \\
\code{>>=} & Right shift and assign \\
\code{>>>=} & Logical right shift and assign \\
\bottomrule
\end{tabular}
\end{table}

\subsection{Other Operators}

\begin{table}[h]
\centering
\begin{tabular}{ll}
\toprule
\textbf{Operator} & \textbf{Description} \\
\midrule
\code{|>} & Pipe (function composition) \\
\code{..} & Range (exclusive end) \\
\code{..=} & Range (inclusive end) \\
\code{<-} & Channel send \\
\code{->} & Function return type \\
\code{=>} & Match arm \\
\code{@} & Decorator (workflow triggers) \\
\bottomrule
\end{tabular}
\end{table}

\section{Delimiters}

\begin{table}[h]
\centering
\begin{tabular}{ll}
\toprule
\textbf{Delimiter} & \textbf{Usage} \\
\midrule
\code{(} \code{)} & Grouping, function calls, tuples \\
\code{[} \code{]} & Array literals, indexing \\
\code{\{} \code{\}} & Blocks, struct literals, maps \\
\code{,} & Separator \\
\code{:} & Type annotations, map entries \\
\code{.} & Member access \\
\code{?} & Optional type marker \\
\bottomrule
\end{tabular}
\end{table}

\section{Whitespace}

Whitespace (spaces, tabs) is used to separate tokens but is otherwise insignificant, except for indentation in multi-line strings.

\begin{lstlisting}
x=1+2         // Valid
x = 1 + 2     // Also valid, preferred for readability
\end{lstlisting}

\section{Lexical Grammar Summary}

\begin{lstlisting}[language=EBNF]
token = keyword
      | identifier
      | integer_literal
      | float_literal
      | string_literal
      | operator
      | delimiter ;

keyword = "agent" | "and" | "as" | "break" | "catch" | "chan"
        | "continue" | "defer" | "else" | "enum" | "false"
        | "fn" | "for" | "from" | "if" | "import" | "in"
        | "match" | "nil" | "not" | "or" | "provider" | "return"
        | "select" | "spawn" | "struct" | "tool" | "true" | "try"
        | "type" | "while" | "workflow" ;

identifier = letter { letter | digit | "_" } ;

operator = "+" | "-" | "*" | "/" | "%"
         | "==" | "!=" | "<" | ">" | "<=" | ">="
         | "and" | "or" | "not" | "&&" | "||" | "!"
         | "&" | "|" | "^" | "~" | "<<" | ">>" | ">>>"
         | "=" | "+=" | "-=" | "*=" | "/="
         | "&=" | "|=" | "^=" | "<<=" | ">>=" | ">>>="
         | "|>" | ".." | "..=" | "<-" | "->" | "=>"
         | "@" ;

delimiter = "(" | ")" | "[" | "]" | "{" | "}"
          | "," | ":" | "." | "?" ;
\end{lstlisting}
